% GENERATED FILE. Edit canonical lessons, then run npm run generate:books.
% canonical-source-hash: fnv1a64:3beb2eaf122ab4d7
% canonical-lessons: LA-C367-cum, LA-C367-ne, LA-C367-quin, LA-C367-quominus, LA-C367-tamquam

\chapter{Since, Lest, But That}
\label{ch:la-since-lest-but-that}
\hlchaptermodality{367}

\begin{chapteropening}
\textbf{By the end of this chapter:} \emph{I can give a reason with cum, an aim with nē, and use quīn, quōminus and tamquam with the subjunctive.}

Complete the last lesson of chapter 367: I can give a reason with cum, an aim with nē, and use quīn, quōminus and tamquam with the subjunctive.
\end{chapteropening}

\section[cum]{cum — when, since, although (with the subjunctive)}

\label{lesson:LA-C367-cum}



\begin{quote}
Before the new one: say the Latin for therefore, then the Latin for for that reason.
\end{quote}

\subsection*{You'll want to know: cum}
\textbf{cum} — \textquotedblleft{}when, since, although (with the subjunctive)\textquotedblright{}.

\begin{quote}
Cum aeger essem, domī mānsī.
\end{quote}

\begin{quote}
\emph{Since I was ill, I stayed at home.}
\end{quote}

As a preposition \emph{cum} takes the ablative: \emph{cum amīcō}, with a friend. As a conjunction with the subjunctive it gives a reason or the scene: \emph{cum essem}, since I was. With the indicative it is simply when.

\begin{cousinweb}
The conjunction \emph{cum} was written \emph{quom} in old Latin (EtymDB; Lewis \& Short): it comes from the stem of the relative \emph{quī} \textquotedblleft{}who, which\textquotedblright{}, and is a different word from the preposition \emph{cum} \textquotedblleft{}with\textquotedblright{}.
\end{cousinweb}

\begin{grammarlens}[title={What you've built}]
One more word for telling, explaining and describing at B1.
\end{grammarlens}

\subsection*{Your turn}
\begin{itemize}
\raggedright
  \item \emph{Say it:} \emph{cum}
  \item \emph{Say it:} \emph{cum}, once more
  \item \emph{Say it:} a sentence of your own with \emph{cum}
  \item \emph{Recall:} say the Latin for therefore, then the Latin for for that reason, then say \emph{cum} again
\end{itemize}

\begin{tcolorbox}[breakable,colback=teal!4,colframe=teal!35!black,title={Before you move on}]
What does cum mean? (\textquotedblleft{}When, since, although (with the subjunctive)\textquotedblright{}.) Say it once more.
\end{tcolorbox}
\section[nē]{nē — so that not, lest; that (after fearing)}

\label{lesson:LA-C367-ne}



\begin{quote}
Before the new one: say the Latin for for that reason, then the Latin for when.
\end{quote}

\subsection*{You'll want to know: nē}
\textbf{nē} — \textquotedblleft{}so that not, lest; that (after fearing)\textquotedblright{}.

\begin{quote}
Portam claudō nē fūr intret.
\end{quote}

\begin{quote}
\emph{I shut the gate so that a thief may not come in.}
\end{quote}

\emph{Nē} + subjunctive gives a negative aim: \emph{ut} says so that, \emph{nē} says so that not. After a verb of fearing, \emph{timeō nē} means I fear that. \emph{Nē … quidem} means not even.

\begin{cousinweb}
\emph{nē} is the old Latin negative; Lewis \& Short compare Sanskrit \emph{na} and Old English \emph{ne}, and it is usually traced to a Proto-Indo-European negative \emph{*ne}. It also lies behind \emph{nōn} and \emph{nihil} and, through Old English, behind English \emph{no}.
\end{cousinweb}

\begin{grammarlens}[title={What you've built}]
One more word for telling, explaining and describing at B1.
\end{grammarlens}

\subsection*{Your turn}
\begin{itemize}
\raggedright
  \item \emph{Say it:} \emph{nē}
  \item \emph{Say it:} \emph{nē}, once more
  \item \emph{Say it:} a sentence of your own with \emph{nē}
  \item \emph{Recall:} say the Latin for for that reason, then the Latin for when, then say \emph{nē} again
\end{itemize}

\begin{tcolorbox}[breakable,colback=teal!4,colframe=teal!35!black,title={Before you move on}]
What does nē mean? (\textquotedblleft{}So that not, lest; that (after fearing)\textquotedblright{}.) Say it once more.
\end{tcolorbox}
\section[quīn]{quīn — but that, that (after a denied doubt); why not}

\label{lesson:LA-C367-quin}



\begin{quote}
Before the new one: say the Latin for when, then the Latin for so that not.
\end{quote}

\subsection*{You'll want to know: quīn}
\textbf{quīn} — \textquotedblleft{}but that, that (after a denied doubt); why not\textquotedblright{}.

\begin{quote}
Nōn dubitō quīn vērum dīcās.
\end{quote}

\begin{quote}
\emph{I do not doubt that you are telling the truth.}
\end{quote}

\emph{Quīn} + subjunctive follows a doubt or a hindrance that is denied: \emph{nōn dubitō quīn}, I do not doubt that. Alone, \emph{quīn} can mean why not.

\begin{cousinweb}
\emph{quīn} is usually taken as the old adverb \emph{quī} \textquotedblleft{}how\textquotedblright{} with the negative \emph{-ne} (Lewis \& Short): how not, why not.
\end{cousinweb}

\begin{grammarlens}[title={What you've built}]
One more word for telling, explaining and describing at B1.
\end{grammarlens}

\subsection*{Your turn}
\begin{itemize}
\raggedright
  \item \emph{Say it:} \emph{quīn}
  \item \emph{Say it:} \emph{quīn}, once more
  \item \emph{Say it:} a sentence of your own with \emph{quīn}
  \item \emph{Recall:} say the Latin for when, then the Latin for so that not, then say \emph{quīn} again
\end{itemize}

\begin{tcolorbox}[breakable,colback=teal!4,colframe=teal!35!black,title={Before you move on}]
What does quīn mean? (\textquotedblleft{}But that, that (after a denied doubt); why not\textquotedblright{}.) Say it once more.
\end{tcolorbox}
\section[quōminus]{quōminus — that … not (after hindering), from doing}

\label{lesson:LA-C367-quominus}



\begin{quote}
Before the new one: say the Latin for so that not, then the Latin for but that.
\end{quote}

\subsection*{You'll want to know: quōminus}
\textbf{quōminus} — \textquotedblleft{}that … not (after hindering), from doing\textquotedblright{}.

\begin{quote}
Nōn recūsō quōminus hoc faciam.
\end{quote}

\begin{quote}
\emph{I do not refuse to do this.}
\end{quote}

\emph{Quōminus} + subjunctive follows verbs of hindering and refusing: \emph{nōn recūsō quōminus}, I do not refuse to.

\begin{cousinweb}
\emph{quōminus} is built from \emph{quō} \textquotedblleft{}by which\textquotedblright{} and \emph{minus} \textquotedblleft{}less\textquotedblright{} (Lewis \& Short): by which the less, so that not.
\end{cousinweb}

\begin{grammarlens}[title={What you've built}]
One more word for telling, explaining and describing at B1.
\end{grammarlens}

\subsection*{Your turn}
\begin{itemize}
\raggedright
  \item \emph{Say it:} \emph{quōminus}
  \item \emph{Say it:} \emph{quōminus}, once more
  \item \emph{Say it:} a sentence of your own with \emph{quōminus}
  \item \emph{Recall:} say the Latin for so that not, then the Latin for but that, then say \emph{quōminus} again
\end{itemize}

\begin{tcolorbox}[breakable,colback=teal!4,colframe=teal!35!black,title={Before you move on}]
What does quōminus mean? (\textquotedblleft{}That … not (after hindering), from doing\textquotedblright{}.) Say it once more.
\end{tcolorbox}
\section[tamquam]{tamquam — as, in the same way as; as if}

\label{lesson:LA-C367-tamquam}



\begin{quote}
Before the new one: say the Latin for but that, then the Latin for that … not.
\end{quote}

\subsection*{You'll want to know: tamquam}
\textbf{tamquam} — \textquotedblleft{}as, in the same way as; as if\textquotedblright{}.

\begin{quote}
Loquitur tamquam rēx sit.
\end{quote}

\begin{quote}
\emph{He speaks as if he were king.}
\end{quote}

\emph{Tamquam} compares: in the same way as. \emph{Tamquam} or \emph{tamquam sī} + subjunctive is as if: \emph{tamquam rēx sit}.

\begin{cousinweb}
\emph{tamquam} is built from \emph{tam} \textquotedblleft{}so\textquotedblright{} and \emph{quam} \textquotedblleft{}as\textquotedblright{} (EtymDB; Lewis \& Short): so as.
\end{cousinweb}

\begin{grammarlens}[title={What you've built}]
One more word for telling, explaining and describing at B1.
\end{grammarlens}

\subsection*{Your turn}
\begin{itemize}
\raggedright
  \item \emph{Say it:} \emph{tamquam}
  \item \emph{Say it:} \emph{tamquam}, once more
  \item \emph{Say it:} a sentence of your own with \emph{tamquam}
  \item \emph{Recall:} say the Latin for but that, then the Latin for that … not, then say \emph{tamquam} again
\end{itemize}

\begin{tcolorbox}[breakable,colback=teal!4,colframe=teal!35!black,title={Before you move on}]
What does tamquam mean? (\textquotedblleft{}As, in the same way as; as if\textquotedblright{}.) Say it once more.
\end{tcolorbox}
